% ==========================================================================
% INFO 4617 Web Data Science -- Week 8: Set up Selenium
% Build: `make` in handouts/, or `latexmk -pdf selenium-setup.tex` here.
% The figures are real captures of Jupyter Notebook 7.6.3, taken and marked
% up by the textbook's tools/shots and copied here by its `sync`; the code in
% them is chapter 8's own (img/IMAGES.md lists each figure's recipe).
% ==========================================================================
\makeatletter\def\input@path{{../common/}{handouts/common/}}\makeatother
\documentclass{handout}

\title[INFO 4617 · Week 8 · Set up Selenium]{Set up Selenium}
\subtitle{Week 8 · before Wednesday's lab · about 15 minutes}
\author{Brian C. Keegan, Ph.D.}
\institute{Department of Information Science, University of Colorado Boulder}
\date{October 5, 2026}

% A numbered figure: [width]{path from the repo root}{alt text}{caption}.
% The picture and its caption share a box, so a page never comes between them.
\newcounter{handoutfig}
\newcommand{\activityfigure}[4][\linewidth]{%
  \par\smallskip\refstepcounter{handoutfig}%
  \noindent\begin{minipage}{\linewidth}%
  {\centering
   \BeginAccSupp{method=pdfstringdef,Alt={#3}}%
   \includegraphics[width=#1]{#2}%
   \EndAccSupp{}\par}%
  \smallskip{\small\raggedright\textbf{Figure \thehandoutfig.} #4\par}%
  \end{minipage}\par}
% A command for a terminal, not a notebook cell: the output's gray bar
\lstnewenvironment{terminal}{\lstset{style=handoutout}}{}

\setlist[enumerate]{leftmargin=*,itemsep=1.5pt,topsep=2pt}
\setlist[itemize]{leftmargin=*,itemsep=1pt,topsep=2pt}

\begin{document}
\maketitle

\begin{frame}{What you do}
  \begin{columns}[T]
    \begin{column}{0.57\textwidth}
      You run four cells from chapter~8 in Jupyter, one at a time, and check
      what each one prints before you run the next. The fourth opens Chrome.
      Then a cell of your own visits a page, prints its title, and closes
      Chrome again.

      \smallskip
      The cells are in chapter~8's section \textbf{Before the First Browser}.
      Copy each one with the \textbf{Copy to Clipboard} button at its top
      right, or open the chapter's companion notebook, which the top of the
      chapter links to.

      \smallskip
      The first run downloads files: the driver, about 10~MB, and, if your
      computer has no Chrome, Chrome for Testing, about 190~MB. Do it on a
      fast connection, not in the minutes before the lab.
    \end{column}
    \begin{column}{0.39\textwidth}
      \begin{block}{Selenium Manager}
        \small A small program that comes with Selenium. It finds your
        Chrome, downloads the driver that matches it, and downloads Chrome
        for Testing if you have no Chrome. It keeps them in
        \texttt{\textasciitilde/.cache/selenium}.
      \end{block}
      \begin{alertblock}{No \texttt{conda activate} needed}
        \small Step~1 finds Selenium Manager through your notebook's own
        Python, however you started Jupyter.
      \end{alertblock}
    \end{column}
  \end{columns}
\end{frame}

\begin{frame}{1 · Settings first}
  \begin{columns}[T]
    \begin{column}{0.52\textwidth}
      Start Jupyter the way you do for every lab, and open a new notebook.
      Paste step~1's cell into it and run it with \texttt{Shift+Enter}. It
      prints two lines (Figure~\ref{fig:settings}):
      \begin{itemize}
        \item ① Selenium's version. Steps~2 and~3 need 4.20 or later.
        \item ② Where Selenium Manager is: in \texttt{webdata}'s own
          folder, under \texttt{bin} on macOS and Linux, or
          \texttt{Scripts} on Windows. The cell looked there and saved the
          path in \texttt{SE\_MANAGER\_PATH}, the setting that
          \texttt{conda activate} makes. If this line says \texttt{inside
          the selenium package}, your Selenium came from \texttt{pip}, and
          that works too.
      \end{itemize}
    \end{column}
    \begin{column}{0.44\textwidth}
      \begin{alertblock}{\texttt{No module named 'selenium'}?}
        \small Install it from a terminal (on Windows, the Anaconda
        Prompt). \texttt{-n webdata} installs into \texttt{webdata} without
        activating it:
\begin{terminal}
conda install -n webdata -c conda-forge selenium
\end{terminal}
        Type \texttt{y} when conda asks. Then, in Jupyter, choose
        \textbf{Kernel} → \textbf{Restart Kernel…} → \textbf{Restart}, and
        run step~1 again. Restart the kernel, not Jupyter: step~1 sets what
        \texttt{conda activate} would have set. For a version older than
        4.20, install \texttt{"selenium>=4.20"} the same way.
      \end{alertblock}
    \end{column}
  \end{columns}
  \activityfigure[0.79\linewidth]{handouts/week-08/img/setup-cell_annotated.pdf}%
    {Screenshot of a Jupyter cell, numbered 1, holding step 1's code. Its
     output has two lines. Marker 1 is beside selenium 4.50.0. Marker 2 is
     beside Selenium Manager: /opt/miniforge3/envs/webdata/bin/selenium-manager.}%
    {Step~1 after it ran: Selenium 4.50.0 ①, and Selenium Manager in
     \texttt{webdata}'s \texttt{bin} folder ②. Two long lines run past the
     cell's right edge, as long lines do in Jupyter.}%
  \label{fig:settings}
  Its last three lines are settings for the steps after it: ignore old
  drivers on your \texttt{PATH}, send no usage statistics, and, commented
  out, a proxy for a network that needs one.
\end{frame}

\begin{frame}{2 · Run Selenium Manager}
  Paste step~2's cell into the next cell and run it. It asks Selenium
  Manager which driver and which browser \texttt{webdriver.Chrome()} will
  use, and shows Selenium Manager's log. The first time, the cell can take a
  minute or more while it downloads.

  The log is pink because it goes to Python's error stream; pink here
  doesn't mean that something failed. Read it from the top
  (Figure~\ref{fig:log}):
  \begin{itemize}
    \item ① \texttt{chrome not found in the system}: this computer has no
      Chrome. If yours has Chrome, the log says \texttt{Detected browser:
      chrome} and its version here, and ② doesn't happen.
    \item ② It downloads Chrome for Testing, a build of Chrome made for
      automation.
    \item ③ It downloads the chromedriver made for that version of Chrome.
    \item ④ Where it keeps the two: in \texttt{\textasciitilde/.cache/selenium},
      or \texttt{\%USERPROFILE\%\textbackslash .cache\textbackslash selenium}
      on Windows. The next run finds them there, and nothing downloads.
  \end{itemize}
  \activityfigure[0.79\linewidth]{handouts/week-08/img/manager-log_annotated.pdf}%
    {Screenshot of a Jupyter cell, numbered 2, holding step 2's code, with
     Selenium Manager's log under it on a pink background. Every log line
     starts with DEBUG. Marker 1 is beside chrome not found in the system.
     Marker 2 is beside Downloading chrome 154.0.8037.92. Marker 3 is
     beside Downloading chromedriver 154.0.8037.92. A brace with marker 4
     spans the last two lines, Driver path and Browser path, both in
     /home/student/.cache/selenium.}%
    {Step~2's first run on a computer with no Chrome and an empty cache
     (October 5, 2026): no Chrome ①, so Selenium Manager downloads Chrome
     for Testing ② and its driver ③, and keeps both in its cache ④.}%
  \label{fig:log}
  Check that no line starts with \texttt{WARNING}. Selenium Manager warns
  instead of stopping, and a warning now can become an error when the
  browser starts.
\end{frame}

\begin{frame}{3 · Check before you start the browser}
  Paste step~3's cell and run it. Start the browser only when its output
  looks like Figure~\ref{fig:check}:
  \begin{itemize}
    \item ① Both paths end in \texttt{(found)}.
    \item ② The last line says \texttt{True}: the driver is the one that
      Selenium Manager matched to your browser, not an old one from your
      \texttt{PATH}.
  \end{itemize}
  \activityfigure[0.79\linewidth]{handouts/week-08/img/check-paths_annotated.pdf}%
    {Screenshot of a Jupyter cell, numbered 3, holding step 3's code. A
     brace with marker 1 spans two output lines: the driver's path and the
     browser's path, each in /home/student/.cache/selenium and each ending
     in (found). Marker 2 is beside Driver chosen by Selenium Manager: True.}%
    {Step~3 after steps~1 and~2: both paths found ①, and the driver is
     Selenium Manager's choice ②.}%
  \label{fig:check}
\end{frame}

\begin{frame}{4 · Start the browser}
  \begin{columns}[T]
    \begin{column}{0.55\textwidth}
      \begin{enumerate}
        \item Paste the cell under \textbf{Starting the Browser} (it holds
          \texttt{driver = webdriver.Chrome()}) and run it. A Chrome window
          opens (Figure~\ref{fig:window}). Leave it alone: your code drives
          it.
        \item In a new cell, visit a page, print its title, and close the
          browser:
\begin{code}
driver.get("https://quotes.toscrape.com/")
print(driver.title)
driver.quit()
\end{code}
      \end{enumerate}
    \end{column}
    \begin{column}{0.41\textwidth}
      \activityfigure{handouts/week-08/img/selenium_browser.png}%
        {Screenshot of the top left of a Chrome for Testing window showing
         xkcd.com. A bar under the address bar says Chrome for Testing
         v154.0.8037.92 is only for automated testing.}%
        {A window that \texttt{webdriver.Chrome()} opened, here on
         xkcd.com. Chrome for Testing says that it is only for automated
         testing; Google Chrome says that it is being controlled by
         automated test software.}%
      \label{fig:window}
    \end{column}
  \end{columns}
  \activityfigure[0.79\linewidth]{handouts/week-08/img/first-chrome_annotated.pdf}%
    {Screenshot of two Jupyter cells. Cell 4 holds driver =
     webdriver.Chrome() and prints Browser: 154.0.8037.92 and Driver:
     154.0.8037.92, which a brace with marker 1 spans. Cell 5 holds the
     three lines that visit quotes.toscrape.com, print its title, and quit;
     marker 2 is beside its output, Quotes to Scrape.}%
    {The browser cell and the visit cell: the browser's and the driver's
     versions ①, and the page's title ②, printed before
     \texttt{driver.quit()} closed the window.}%
  \label{fig:first}
  The two versions ① should match up to the first dot. End every session
  with \texttt{driver.quit()}: each driver is a whole browser, and one you
  forget keeps running and using memory.
\end{frame}

\begin{frame}{If something goes wrong}
  \textbf{\texttt{NoSuchDriverException: Message: Unable to obtain driver
  for chrome}} at the bottom of a long red traceback, and
  \textbf{\texttt{Unable to obtain working Selenium Manager binary}} partway
  up it (Figure~\ref{fig:nomanager}). You started the browser before
  step~1, in a Jupyter that started without \texttt{webdata} active or
  before you installed Selenium, so Selenium couldn't find Selenium Manager.
  Run step~1, then the browser cell again.
  \activityfigure[0.74\linewidth]{handouts/week-08/img/no-manager.png}%
    {Two strips of a red Jupyter traceback, one above the other. The top
     strip, labeled Partway down the traceback, says WebDriverException:
     Message: Unable to obtain working Selenium Manager binary, then a path
     inside the selenium package, then The above exception was the direct
     cause of the following exception. The bottom strip, labeled At the
     bottom, shows the last lines of Selenium's code and then
     NoSuchDriverException: Message: Unable to obtain driver for chrome,
     with a link to Selenium's documentation.}%
    {The browser cell run without step~1, in a Jupyter started without
     \texttt{SE\_MANAGER\_PATH}: the cause, partway down the traceback, and
     the message at its bottom.}%
  \label{fig:nomanager}

  \begin{itemize}
    \item \textbf{\texttt{This version of ChromeDriver only supports Chrome
      version} \textit{N}}: an old chromedriver on your \texttt{PATH} won. Step~1
      tells Selenium Manager to skip it: run steps~1 to~3, then the browser
      cell.
    \item \textbf{The downloads in step~2 fail}: your network may block
      Google's download servers. Run steps~1 and~2 once on another network,
      such as your home Wi-Fi; after that, Selenium Manager works from its
      cache.
    \item \textbf{Ubuntu 24.04 or later, and \texttt{session not created:
      Chrome instance exited}}: Ubuntu doesn't let Chrome for Testing use
      its sandbox. Install Google Chrome from
      \url{https://www.google.com/chrome/}, and run steps~2 to~4 again;
      Selenium Manager then finds it. Don't add \texttt{-{}-no-sandbox}.
    \item \textbf{Something else}: bring the whole message, from its first
      red line to its last, to Wednesday's lab. Chapter~8's ``When Selenium
      Manager Fails'' lists more causes.
  \end{itemize}
\end{frame}

\end{document}
